% ID: SEL-D0015-BB-E03
% ESTADO: validada. S016. Revisión matemática, didáctica y visual.
\textbf{a) Paso por el punto y condición de extremo.}\par\nopagebreak[4]
Un extremo relativo en $(2,36)$ exige $f(2)=36$ y $f'(2)=0$. Como $f'=3ax^2+b$,
\[
8a+2b+4=36,\qquad12a+b=0.
\]
La primera equivale a $4a+b=16$. Restando la segunda, $-8a=16$, de donde
\[
\boxed{a=-2,\qquad b=24.}
\]
Se comprueba que $f(2)=36$. Además, $f''(2)=12a=-24<0$, por lo que es efectivamente un máximo relativo, no sólo un punto estacionario.

\textbf{b) Monotonía para a=4 y b=-3.}\par\nopagebreak[4]
Estos parámetros son distintos de los obtenidos en a. Ahora $f(x)=4x^3-3x+4$ y
\[
f'(x)=12x^2-3=3(2x-1)(2x+1).
\]
La derivada es positiva para $x<-1/2$ y $x>1/2$, y negativa entre ambas raíces. Así, crece en $(-\infty,-1/2)$ y $(1/2,+\infty)$ y decrece en $(-1/2,1/2)$.
\[
\boxed{\text{Máximo relativo }(-1/2,5);\qquad\text{mínimo relativo }(1/2,3).}
\]
Los cambios de signo justifican ambos extremos; no son extremos absolutos de la función cúbica en $\mathbb R$.

\textbf{c) Primitiva con condición inicial.}\par\nopagebreak[4]
Para los mismos parámetros de b,
\[
F(x)=x^4-\frac32x^2+4x+K.
\]
La condición $F(2)=10$ da $16-6+8+K=10$, es decir, $K=-8$:
\[
\boxed{F(x)=x^4-\frac32x^2+4x-8.}
\]
Derivando se recupera $f$ y sustituyendo $2$ se obtiene $10$. Se conserva el literal tipográfico del enunciado en a sin modificar su puntuación ni el problema matemático.
